\documentclass[12pt,a4paper]{article}
\usepackage[utf8]{inputenc}
\usepackage[T1]{fontenc}
\usepackage{hyperref}
\usepackage{amsmath}
\usepackage{amsfonts}
\usepackage{amssymb}
\usepackage{listings}
\usepackage{xcolor}

\lstset{
    language=Haskell,
    basicstyle=\ttfamily\small,
    keywordstyle=\color{blue},
    commentstyle=\color{gray},
    stringstyle=\color{red},
    breaklines=true,
    showstringspaces=false,
    frame=single,
    morekeywords={deriving,newtype,data,instance,where,type,kind}
}

\title{Outline: Minimal Functors, Applicatives, and Monads in Haskell}
\author{Documentation}
\date{\today}

\begin{document}

\maketitle

\section{Introduction}
\begin{itemize}
    \item \textbf{The Journey}: Constructing minimal instances (like \texttt{Proxy} and \texttt{Identity}) to strip away domain-specific noise and reveal the pure mechanics of Haskell's core typeclasses.
    \item \textbf{The Narrative}: Noting that while these core concepts are foundational to modern Haskell, this specific teaching narrative---starting with absolute minimalism to ``prove'' the forced hand of parametricity---is unique, synthesizing ideas from papers like Wadler's ``Theorems for Free!'' and books like Maguire's ``Thinking with Types''.
    \item \textbf{The Scope}: Establishing that we are focusing specifically on Endofunctors operating within the category of Haskell types (from \texttt{Hask} to \texttt{Hask}).
    \begin{itemize}
        \item \textbf{Note on \texttt{Hask}}: Mention that \texttt{Hask} is technically not a strict category due to non-terminating programs (bottom, \texttt{\_|\_}). Include a link to the famous paper \href{https://www.cse.chalmers.se/~nad/publications/danielsson-et-al-popl2006.pdf}{``Fast and Loose Reasoning is Morally Correct''}.
    \end{itemize}
    \item \textbf{The Power of Parametricity}: Exploring how the inability to inspect types at runtime forces implementations to be mathematically unique—the ``theorems for free'' philosophy.
    \begin{itemize}
        \item \textit{Counterexamples}: Add a very brief note acknowledging rare counterexamples where multiple valid implementations might exist (e.g., variations in ordering/traversal), but confirm they are atypical for minimal types.
    \end{itemize}
\end{itemize}

\section{The Foundations of Functors, Applicatives, and Monads}

\subsection{What is a Functor?}
\begin{itemize}
    \item \textbf{The Original Definition}: Formalizing Functors in Category Theory (citing Saunders Mac Lane) as a mapping $F: C \to D$ consisting of:
    \begin{itemize}
        \item \textbf{Object Mapping}: For every object $X$ in $C$, an object $F(X)$ in $D$.
        \item \textbf{Morphism Mapping}: For every morphism $f: X \to Y$ in $C$, a morphism $F(f): F(X) \to F(Y)$ in $D$, preserving identities and composition.
    \end{itemize}
    \item \textbf{Programming Ubiquity}: Noting that functors are widely used in programming languages (even non-functional ones) without developers always knowing it.
    \item \textbf{Haskell's Category Theory Context}: Highlighting that a Haskell \texttt{Functor} is an \textit{Endofunctor} mapping from \texttt{Hask} to \texttt{Hask}.
    \begin{itemize}
        \item \textbf{Examples of Non-Haskell Categorical Functors}:
        \begin{itemize}
            \item \textbf{Non-Endofunctor}: A functor mapping from a completely different category (e.g., the category of Sets) to \texttt{Hask}.
            \item \textbf{Non-parametric Functor}: A categorical functor that inspects types (e.g., mapping \texttt{Int} to \texttt{String} and \texttt{Bool} to \texttt{Double}). This is strictly forbidden in Haskell due to parametricity.
            \item \textbf{Restricted Functor}: A generic categorical functor might only apply to a \textit{subset} of objects (types). A Haskell \texttt{Functor} mathematically must unconditionally apply to \textit{all} types.
        \end{itemize}
    \end{itemize}
    \item \textbf{Taxonomy}: Differences between \texttt{type}, \texttt{newtype}, and \texttt{data}.
\end{itemize}

\subsection{The Constraint of Parametricity}
\begin{itemize}
    \item \textbf{The Concept}: How \texttt{fmap} generically handling \textit{all} types forces it to preserve structure.
    \item \textbf{Theorems for Free}: Linking to Wadler's paper.
\end{itemize}

\subsection{\texttt{fmap} (and \texttt{<\$>}) and the Functor Laws}
\begin{itemize}
    \item \textbf{The Signature}: \texttt{(a -> b) - > f a -> f b}. Lifting function application into a context.
    \item \textbf{The Laws}: Identity and Composition. 
    \item \textbf{Enforcement}: Use of \texttt{QuickCheck} for automated verification.
\end{itemize}

\subsection{The Applicative Functor}
\begin{itemize}
    \item \textbf{The Context}: Functor equipped with \texttt{pure} and \texttt{apply} (\texttt{<*>}).
    \item \textbf{The Laws}: Briefly listing Applicative laws.
\end{itemize}

\subsection{The Monad}
\begin{itemize}
    \item \textbf{Three Equivalent Paths}: \texttt{bind}, \texttt{join} (\texttt{mu}), and \texttt{kleisli} composition (\texttt{>=>}).
\end{itemize}

\section{The Minimal Functor/Monad (Proxy)}
\textit{(Zero computational data, Zero contextual data)}

\subsection{The Smallest Valid Candidate}
\begin{itemize}
    \item \textbf{The Definition}: \texttt{data MinF a = Val}.
    \item \textbf{Haskell Equivalents}: \texttt{Proxy} or \texttt{Const ()}.
\end{itemize}

\subsection{A Singular Functor Implementation}
\begin{itemize}
    \item \textbf{The Code}: \texttt{fmap \_ Val = Val}.
    \item \textbf{The ``Why''}: The type signature forces our hand.
\end{itemize}

\subsection{Upgrading to Applicative}
\begin{itemize}
    \item \textbf{The Code}: \texttt{pure \_ = Val} and \texttt{Val <*> Val = Val}.
\end{itemize}

\subsection{Upgrading to Monad}
\begin{itemize}
    \item \textbf{Via \texttt{join}}: \texttt{join Val = Val}.
    \item \textbf{Via \texttt{bind}}: \texttt{Val >>= \_ = Val}.
\end{itemize}

\subsection{Let the Compiler Do the Work (\texttt{deriving})}
\begin{itemize}
    \item Using \texttt{DeriveFunctor}.
\end{itemize}

\section{The Minimal Applicative Functor (Const)}
\textit{(Zero computational data, Some contextual data \texttt{r}). Requires \texttt{Monoid r}.}

\subsection{The Definition of \texttt{Const}}
\begin{itemize}
    \item \texttt{newtype Const r a = Const r}.
\end{itemize}

\subsection{A Singular Functor Implementation}
\begin{itemize}
    \item \texttt{fmap \_ (Const r) = Const r}.
\end{itemize}

\subsection{The Applicative Twist (Necessary and Sufficient)}
\begin{itemize}
    \item Why \texttt{Monoid r} is strictly necessary for \texttt{pure} (identity) and \texttt{<*>} (combination).
\end{itemize}

\subsection{Why Not a Monad?}
\begin{itemize}
    \item Failing the Left Identity law.
\end{itemize}

\section{The Minimal Synchronous Monad (Identity)}
\textit{(One computational data, Zero contextual data)}

\subsection{The Definition}
\begin{itemize}
    \item \texttt{data IdF a = IdVal a}.
\end{itemize}

\subsection{A Singular Functor Implementation}
\begin{itemize}
    \item \texttt{fmap f (IdVal x) = IdVal (f x)}.
\end{itemize}

\subsection{Upgrading to Applicative}
\begin{itemize}
    \item \texttt{pure x = IdVal x} and \texttt{IdVal f <*> IdVal x = IdVal (f x)}.
\end{itemize}

\subsection{Upgrading to Monad}
\begin{itemize}
    \item \texttt{join (IdVal (IdVal x)) = IdVal x}.
    \item \texttt{IdVal x >>= f = f x}.
\end{itemize}

\section{The Algebra of Functors}
\begin{itemize}
    \item \textbf{Foundational Blocks}: Treating \texttt{Proxy} as the number $1$ and \texttt{Identity} as the variable $X$ in a type-level polynomial.
    \item \textbf{Combinations}: How the concepts of ``OR'' (Sums) and ``AND'' (Products) allow us to build complex ADTs from minimal primitives.
\end{itemize}

\section{Functors out of Proxy and Identity}
\begin{itemize}
    \item \textbf{The Sum (Maybe)}: Deriving \texttt{Maybe} as the sum of Zero and One ($1 + X$), representing the minimal monad of failure.
    \item \textbf{The Product (Identity)}: Proving that the product of Zero and One ($1 \times X$) collapses back into \texttt{Identity}, showing that \texttt{Proxy} acts as the multiplicative identity.
    \item \textbf{Extensions}: Generalizing the product to \texttt{Writer} by replacing the unit with a Monoidal constant \texttt{r}.
\end{itemize}

\section{Functors entirely out of Proxy}
\begin{itemize}
    \item \textbf{Fundamental Rule}: Reinforcing that \texttt{Proxy} represents a single computational state with no attached data.
    \item \textbf{Proxy + Proxy}: Showing how summing two empty states ($1 + 1 = 2$) creates the \texttt{Const Bool} functor.
    \item \textbf{Proxy * Proxy}: Proving that a tuple of units ($1 \times 1 = 1$) remains isomorphic to a single unit.
\end{itemize}

\section{Recursion and Fixed Points}
\begin{itemize}
    \item \textbf{Recursive Sums/Products}: Building the \texttt{List} type algebraically as $L(X) = 1 + X \times L(X)$.
    \item \textbf{Shape Equations}: Understanding how different polynomial formulas define the physical layout of data structures like Trees and Streams.
    \item \textbf{Alternative Shapes}: Brief comparison of shape formulas for different recursive types.
\end{itemize}

\section{Conclusion: The Tale of Three Minimals}
\begin{itemize}
    \item \textbf{Comparison}: Zero (\texttt{MinF}), Accumulation (\texttt{Const}), and One (\texttt{Identity}).
\end{itemize}

\section{Annex \& Bibliography}
\begin{itemize}
    \item Proofs of Monad equivalence and unique Functor identity.
\end{itemize}

\end{document}
